%% ma: L12-B17-0011
%% lop: 12
%% chuong: 5
%% bai: 17
%% dang: 12.17.1
%% loai: TLN
%% muc: VD
%% dap_an: "240"
%% nguon: "(NBV)-ĐỀ SỐ 19-MỖI NGÀY 1 ĐỀ THI 2025.pdf (D19, MỖI NGÀY 1 ĐỀ - Đề số 19), Phần III Câu 6 (Sở Ninh Bình 2025)"
%% kien_thuc: "Bài 15 (phương trình tham số của đường thẳng), Bài 17 (mặt cầu, giao của đường thẳng và mặt cầu); chuyển động thẳng đều"
%% trang_thai: chinh-thuc
%% ghi_chu: "Đáp án gốc 240 đúng (giải độc lập: |OB|² = 144 ⇔ t² - 8t + 12 = 0 ⇔ t = 2 hoặc 6; thời gian 4 giờ = 240 phút). Cùng ý tưởng với L12-B17-0001 và L12-B17-0008 (vật thể bay và vùng theo dõi hình cầu) nên nhập vào dạng đã duyệt 12.17.1. Bỏ ảnh minh họa tháp kiểm soát và khinh khí cầu (ảnh trang trí)."
\begin{ex}
Trong không gian $Oxyz$, một khinh khí cầu ở tọa độ $A(-16;-10;10)$ bắt đầu bay với vectơ vận tốc không đổi $\vec v=(4;3;-1)$ (đơn vị vận tốc là km/h) và dự kiến bay trong thời gian $10$ giờ. Biết trạm kiểm soát không lưu đặt ở vị trí gốc tọa độ $O$ kiểm soát được các vật thể cách trạm một khoảng tối đa bằng $12$ km. Thời gian kể từ khi trạm kiểm soát không lưu phát hiện ra khinh khí cầu đến khi khinh khí cầu ra khỏi vùng kiểm soát là bao nhiêu phút?
\shortans{240}
\loigiai{Sau $t$ giờ ($0\le t\le10$) khinh khí cầu ở $B(-16+4t;-10+3t;10-t)$. Trạm kiểm soát được khi $OB\le12$:
$$(-16+4t)^2+(-10+3t)^2+(10-t)^2\le144\Leftrightarrow26t^2-208t+312\le0\Leftrightarrow t^2-8t+12\le0\Leftrightarrow2\le t\le6.$$
Cả hai mốc đều nằm trong thời gian bay dự kiến $[0;10]$. Vậy khinh khí cầu ở trong vùng kiểm soát $6-2=4$ giờ, tức $240$ phút.}
\end{ex}
